\section{Introduction}
\label{sec:introduction}

Retrieval-augmented generation (RAG) grounds language-model answers in external evidence, but multi-hop question answering stresses this design. A question may require linking entities, comparing facts, or combining evidence from multiple documents. If retrieval returns only part of the supporting evidence, the reader may still fail even when the answer is straightforward from the full evidence set. Supporting-title completeness, checked against supporting paragraph text when the processed artifacts expose it, is therefore a natural diagnostic for multi-hop RAG.

Prior work addresses multi-hop evidence acquisition with more adaptive retrieval and reasoning. IRCoT-style methods interleave reasoning traces with retrieval, and agentic or self-reflective RAG systems can decide when to retrieve, critique, or revise outputs. We ask a narrower empirical question: under a fixed-depth, training-free protocol, when does HyDE-style generative expansion improve complete supporting-title and paragraph-text retrieval without an open-ended retrieval loop?

Building on HyDE rather than proposing a new retriever architecture, we examine document-like generative query expansion in this fixed-depth protocol. The evaluated pipeline generates a concise hypothetical supporting passage, retrieves with the question plus that passage as the dense query, and then applies an evidence-grounded short-answer reader. Unlike direct query rewriting, the generated text is document-like and may expose likely entities, relations, and answer types for dense retrieval.

The main empirical pattern is evidence-driven. On HotpotQA \texttt{test\_subsampled}, HyDE-style RAG raises all-support hit@10 to 0.8680 and mean supporting-title recall@10 to 0.9280, with reader EM/F1 of 0.6840/0.8105. These values exceed the implemented lightweight retrieval baselines and diagnostic controls in our fixed local-corpus protocol, including dense, BM25, hybrid, rule-based multi-query, rule-based iterative retrieval, and matched-call query reformulation. We do not treat these author-implemented controls as faithful reproductions of all neighboring published systems. A secondary guarded canonicalization diagnostic produces only a small, statistically non-significant answer-format change, so we treat it as exploratory post-processing rather than as a primary accuracy contribution.

This controlled study makes three empirical contributions. First, it shows that HyDE-style document-like expansion can improve complete supporting-title retrieval in a fixed-depth multi-hop RAG protocol, with matching paragraph-text audit results and answer gains reported for the fixed reader used in the main evaluation. Second, it isolates competing explanations through matched-call, query-composition, paired, supporting-title and paragraph-text, answer-string-overlap, strict equal-budget query-form, and top-$k$ sensitivity diagnostics; these controls reduce explanations based only on an extra LLM call, retained question text, ordinary direct rewriting, generated length, title-level scoring artifacts, or a single cutoff at rank 10. Third, it characterizes boundary conditions through answer-overlap, gold-content scrubbing, stronger-encoder, corpus-scale leakage and hard-negative audits, and targeted reader checks, showing that answer-like generated content, high-information keyword expansion, encoder strength, corpus scale and difficulty, sentence-level annotation limits, and reader choice affect how far the mechanism claim should generalize. The resulting claim is intentionally retrieval-stage and diagnostic: document-like expansion improves complete supporting-title and paragraph-text acquisition over direct reformulation and question decomposition under the evaluated protocol, while its advantage over keyword/entity expansion is dataset-dependent. We additionally report guarded canonicalization only as a secondary answer-format risk analysis.
